\section{Dodecaphasic polarization, lattice duality, and theta reciprocity}
\label{sec:leech-dualidad}

The exceptional incidence and the polarization of the dodecaphasic record
arise from two realizations of the same enriched state. The first
preserves the boundary relations, the glue code, and the marked
origin determining the lattice; the second lets a
representation act on the twelve integer coordinates already obtained by the
readers of the Topological Phase Kernel. Their composition allows us to study
a metric deformation whose parameter changes the relative lengths
of twelve planes while simultaneously preserving the total volume and an
explicit ternary symmetry.

The causal starting point remains the APP--TRIT--TPK
construction of the enriched state and the joint continuum.
The record \(K\) is received after the dodecaphasic record and its
incidence and orientation operations. The representation used
acts on this output and preserves its twelve values as already
determined data. The lattice realization is received with its bilinear form,
its glue classes, and its marked vector. These structures, developed
respectively in \cite{hmtX,HMT-IV}, are retained when composed.
The collection constructed below is a geometric realization
of these data with a common twelve-position chart.

\subsection{Polarization of the record and marking of the twelve planes}
\label{subsec:leech-polarizacion}

The dodecaphasic record and its sum are
\begin{equation}
 \begin{split}
 K={}&(234,543,140,729,659,824,\\
     &\hspace{13mm}621,58,914,794,146,601),\\
 \sum_{j=1}^{12}K_j={}&6263.
 \end{split}
 \label{eq:leech-K}
\end{equation}
Let \(A_5\) be the group of even permutations of five elements and let
\(C_5=\langle(1\,2\,3\,4\,5)\rangle\). Its twelve left cosets
determine a permutation representation
\(\varrho:A_5\to O(\mathbb R^{12})\). The marked chart is fixed by
the following representatives, written as lists of images:
\[
\begin{array}{c|c@{\qquad}c|c}
j&r_j&j&r_j\\ \hline
1&(4,3,5,2,1)&7&(5,4,3,2,1)\\
2&(4,2,1,3,5)&8&(1,3,4,2,5)\\
3&(1,2,4,5,3)&9&(4,2,3,5,1)\\
4&(2,5,3,4,1)&10&(3,2,5,4,1)\\
5&(4,3,1,5,2)&11&(4,5,2,3,1)\\
6&(5,1,2,3,4)&12&(5,3,2,4,1).
\end{array}
\]
Specifically, \(\varrho(a)e_j=e_\ell\) if
\(ar_jC_5=r_\ell C_5\). The three-dimensional character \(\chi_3\) takes
the values
\[
\begin{array}{c|ccccc}
\text{class}&1&2^2&3&5_a&5_b\\ \hline
\chi_3&3&-1&0&(1+\sqrt5)/2&(1-\sqrt5)/2.
\end{array}
\]
The class \(5_a\) contains the stated generator of \(C_5\), and \(5_b\)
contains its square. This choice distinguishes the two conjugate
three-dimensional characters. If \(I_{12}\) is the identity and
\(\mathbf1\) is the vector of twelve ones, define
\begin{equation}
 P_3=\frac3{60}\sum_{a\in A_5}\chi_3(a^{-1})\varrho(a),
 \qquad
 P_{11}=I_{12}-\frac1{12}\mathbf1\mathbf1^{\mathsf T}.
 \label{eq:leech-projectors}
\end{equation}
The field \(\mathbb Q(\sqrt5)\) appears here in the realization of the
representation acting on the generated record.

\begin{lemma}[Zero-sum polarization]
\label{lem:leech-u}
The operators in \eqref{eq:leech-projectors} satisfy
\[
 P_3=P_3^{\mathsf T}=P_3^2,\qquad
 \operatorname{rank}P_3=3,\qquad P_3\mathbf1=0,\qquad
 P_3P_{11}=P_{11}P_3=P_3.
\]
Consequently,
\begin{equation}
 u=P_3P_{11}K=P_3K,\qquad
 \sum_{j=1}^{12}u_j=0,\qquad
 \|u\|^2=\frac{6638585+2275584\sqrt5}{20}>0.
 \label{eq:leech-u-norm}
\end{equation}
\end{lemma}
\begin{proof}
The character sum is the isotypic projector of the representation.
The reality of the character and the substitution \(a\mapsto a^{-1}\)
give its symmetry. The multiplicity of the three-dimensional character in the
representation of \(A_5/C_5\) is the dimension of its
\(C_5\)-invariants, which equals
\[
 \frac15\left(
 3+2\frac{1+\sqrt5}{2}+2\frac{1-\sqrt5}{2}
 \right)=1.
\]
The projector therefore has rank three. Its orthogonality to the trivial
character implies \(P_3\mathbf1=0\) and the identities with \(P_{11}\).
Applying the finite sum to the record
\eqref{eq:leech-K} produces the twelve coordinates
\begin{equation}
\begin{aligned}
u={}&(-495/4-233\sqrt5/10,\quad403/4+169\sqrt5/10,\\
&-403/4-169\sqrt5/10,\quad495/4+233\sqrt5/10,\\
&601/4+1031\sqrt5/10,\quad203/4+211\sqrt5/5,\\
&-203/4-211\sqrt5/5,\quad-601/4-1031\sqrt5/10,\\
&313/4+569\sqrt5/10,\quad162+219\sqrt5/20,\\
&-162-219\sqrt5/20,\quad-313/4-569\sqrt5/10).
\end{aligned}
\label{eq:leech-u-coordinates}
\end{equation}
Summing them and squaring them in \(\mathbb Q(\sqrt5)\) gives
\eqref{eq:leech-u-norm}. Equivalently,
\[
 \|u\|^2=K^{\mathsf T}P_3K
 =\frac1{20}\sum_{a\in A_5}
          \chi_3(a^{-1})K^{\mathsf T}\varrho(a)K.
\]
The positive coefficients of its evaluation prove that \(u\ne0\).
\end{proof}

The zero sum has a concrete geometric function: expansion of
some planes can compensate for contraction of the others in the
total determinant. Non-vanishing of \(u\), by contrast, ensures that this
compensation contains an actual deformation.

The ambient Euclidean space of the lattice realization is
\begin{equation}
 E=\bigoplus_{j=1}^{12}E_j,\qquad
 B=\operatorname{diag}(B_2,\ldots,B_2),\qquad
 B_2=\begin{pmatrix}2&-1\\-1&2\end{pmatrix}.
 \label{eq:leech-ambient}
\end{equation}
Each \(E_j\) is the real plane spanned by the root lattice
\(A_2\), expressed in its simple-root basis. The orthogonal
decomposition belongs to the ambient space. The Leech lattice is
obtained by gluing and passage to a neighbour, operations relating
the twelve factors.

Let \(C_W\subset\mathbb F_3^{12}\) be the six-dimensional ternary code
of the Paley--Witt construction. Identifying the discriminant of each
\(A_2\) factor with \(\mathbb F_3\), the gluing is
\[
 N=N(C_W)=
 \bigcup_{c\in C_W}\bigl(A_2^{12}+\phi(c)\bigr),
\]
where \(\phi(c)\) chooses representatives of the discriminant classes.
Self-orthogonality and self-duality of the code transmit evenness and
integrality to \(N\); its determinant is
\[
 \det N=\frac{3^{12}}{(3^6)^2}=1.
\]
The exceptional construction fixes the marked origin \(j=1\) and the
vector
\begin{equation}
 v_\alpha=4\rho_1+\rho_2+\cdots+\rho_{12},
 \qquad \rho_j=(1,1)_j,\qquad
 \|v_\alpha\|_B^2=54.
 \label{eq:leech-marked-vector}
\end{equation}
Using it, we form
\begin{equation}
\begin{split}
 N_{\alpha,0}
  &=\{x\in N:\langle x,v_\alpha\rangle_B\equiv0\pmod3\},\\
 \Lambda&=N_{\alpha,0}+\mathbb Z\,v_\alpha/3.
\end{split}
\label{eq:leech-neighbour}
\end{equation}
The exceptional realization establishes that \(\Lambda\) is a positive,
even, unimodular, rootless lattice of rank \(24\), so it is
isometric to the Leech lattice \cite{HMT-IV}. In particular,
\(\Lambda^*=\Lambda\), with dual taken with respect to \(B\).
These properties and the membership \(v_\alpha/3\in\Lambda\) are the
precise lattice hypotheses used in the following theorems.

\subsection{Ternary isometry and stability of the neighbour}
\label{subsec:leech-c3}

We define the two-dimensional matrices
\[
 C=\begin{pmatrix}0&-1\\1&-1\end{pmatrix},
 \qquad
 R=\begin{pmatrix}0&1\\1&0\end{pmatrix},
 \qquad
 \varpi=\frac13\begin{pmatrix}2\\1\end{pmatrix}.
\]
Direct multiplication verifies
\[
 C^3=I_2,\quad C^2+C+I_2=0,\quad
 C^{\mathsf T}B_2C=B_2,\quad
 RCR=C^{-1},\quad R^{\mathsf T}B_2R=B_2.
\]
Moreover, \(C\varpi-\varpi=(-1,0)^{\mathsf T}\), so \(C\)
acts trivially on \(A_2^*/A_2\). The reflection \(R\) reverses the sign
of the discriminant class. The code \(C_W\) contains the full-support
word
\[
 c_*=(1,1,1,1,1,1,2,1,1,1,1,1).
\]
In the presentation \(C_W=\{(q,qA_W):q\in\mathbb F_3^6\}\), this
word arises from \(q_*=(1,1,1,1,1,1)\) and
\(q_*A_W=(2,1,1,1,1,1)\).

For each plane we fix the frame
\begin{equation}
 Q_j=\begin{cases}
 R,&(c_*)_j=1,\\
 I_2,&(c_*)_j=2,
 \end{cases}
 \qquad
 g_c=\bigoplus_{j=1}^{12}Q_jCQ_j^{-1}.
 \label{eq:leech-c3}
\end{equation}
The symbols \(Q_j\) denote plane frames; the symbol \(P_3\)
is reserved for the projector onto the three-dimensional sector.

\begin{proposition}[Order-three symmetry of the marked lattice]
\label{prop:leech-c3}
The transformation \(g_c\) is an order-three isometry of \(E\),
has zero as its only fixed vector, and satisfies
\[
 g_cN=N,\qquad g_cN_{\alpha,0}=N_{\alpha,0},
 \qquad g_c\Lambda=\Lambda.
\]
\end{proposition}
\begin{proof}
Each block is \(C\) or \(C^{-1}\), whose minimal polynomial is
\(X^2+X+1\). This proves the order and the claim about fixed vectors,
as well as preservation of \(B\). Each block acts trivially on
the discriminant and therefore preserves all glue classes of
\(N\).

Write \(v=v_\alpha=\sum_j a_j\rho_j\), with \(a_1=4\) and
\(a_j=1\) for \(j>1\). All \(a_j\) are congruent to one
modulo three. The identities
\[
 (C-I_2)\rho=(-2,-1)^{\mathsf T},
 \qquad (C^{-1}-I_2)\rho=(-1,-2)^{\mathsf T}
\]
show that
\[
 y=\frac{g_cv-v}{3},\qquad z=\frac{g_c^{-1}v-v}{3}
\]
have discriminant classes \(c_*\) and \(-c_*\), respectively.
Both belong to \(C_W\), so \(y,z\in N\). In each factor,
\[
 \left\langle
 \frac{(C^{\pm1}-I_2)a_j\rho}{3},a_j\rho
 \right\rangle_{B_2}=-a_j^2.
\]
Summing gives \(\langle y,v\rangle_B=\langle z,v\rangle_B=-27\);
hence \(y,z\in N_{\alpha,0}\).
If \(x\in N_{\alpha,0}\), then
\[
 \langle g_cx,v\rangle_B
  =\langle x,g_c^{-1}v\rangle_B
  =\langle x,v\rangle_B+3\langle x,z\rangle_B
  \equiv0\pmod3.
\]
The same operation with \(g_c^{-1}\) proves the equality
\(g_cN_{\alpha,0}=N_{\alpha,0}\). Finally,
\(g_c(v/3)=v/3+y\) preserves the class generating
\(\Lambda=N_{\alpha,0}+\mathbb Z(v/3)\).
\end{proof}

The ternary symmetry therefore acts on the complete glued
lattice. Its preservation uses the codeword and
the congruence of the marked vector, in addition to the rotation of each
plane. This precision allows the symmetry to be transported to the flow without
losing the arithmetic information that made its construction possible.

\subsection{Anisotropic flow and covolume preservation}
\label{subsec:leech-flow}

\begin{definition}[Deformation determined by polarization]
For a dimensionless parameter \(s\in\mathbb R\), define
\begin{equation}
 H_u=\bigoplus_{j=1}^{12}u_jI_2,\qquad
 D_s=\exp(sH_u)=\bigoplus_{j=1}^{12}e^{su_j}I_2,\qquad
 \Lambda_s=D_s\Lambda.
 \label{eq:leech-flow}
\end{equation}
The coordinate \(u_j\) is assigned to the plane \(E_j\) of the same marked
chart.
\end{definition}

The positive factor \(e^{su_j}\) preserves the orientation of each
plane. The twelve labels transmit the correspondence between
polarization and the lattice realization. This correspondence
is part of the starting structure; a permutation of the
labels is handled by transporting it simultaneously in both objects.

\begin{theorem}[Covolume and symmetry of the flow]
\label{thm:leech-flow}
For all \(s,t\in\mathbb R\), we have
\[
\begin{gathered}
 D_{s+t}=D_sD_t,\qquad D_0=I,\qquad D_s^{-1}=D_{-s},\\
 D_s^\dagger=D_s,\qquad D_sg_c=g_cD_s,\qquad \det D_s=1,
\end{gathered}
\]
where \(A^\dagger=B^{-1}A^{\mathsf T}B\).
Each \(\Lambda_s\) is a discrete lattice of rank \(24\) and covolume
one, preserved by the order-three isometry \(g_c\).
\end{theorem}
\begin{proof}
The group identities hold in each exponential block.
The scalar blocks are self-adjoint for \(B_2\), and commute
with the blocks of \(g_c\). The zero sum in
\eqref{eq:leech-u-norm} gives
\[
 \det D_s=\prod_{j=1}^{12}e^{2su_j}
 =\exp\left(2s\sum_{j=1}^{12}u_j\right)=1.
\]
The image of a lattice under an invertible operator is a lattice of the
same rank; its covolume is multiplied by the absolute value of the
determinant. The covolume of \(\Lambda\) is one, so that
of \(\Lambda_s\) is also one. Commutation and
Proposition~\ref{prop:leech-c3} give
\[
 g_c\Lambda_s=g_cD_s\Lambda=D_sg_c\Lambda=D_s\Lambda.
\]
The order and the fixed-point space of \(g_c\) remain
unchanged.
\end{proof}

\begin{remark}[Parameter and preserved symmetry]
The direction \(u\) is determined by the record; \(s\) ranges over
a collection of realizations. Replacing \(u\) by \(u/\|u\|\)
reparametrizes the same collection. A physical interpretation
selecting a particular value of \(s\) adds its realization rule.
The theorem preserves the explicit symmetry \(g_c\); transporting
other automorphisms requires their corresponding relations with
\(H_u\).
\end{remark}

\subsection{Euclidean dual and actual deformation}
\label{subsec:leech-dual}

\begin{definition}[Euclidean dual]
For a lattice \(L\subset E\), its dual with respect to \(B\) is
\[
 L^*=\{y\in E:\langle y,x\rangle_B\in\mathbb Z
                         \text{ for every }x\in L\}.
\]
\end{definition}

\begin{theorem}[Dual reciprocity]
\label{thm:leech-dual}
The collection \eqref{eq:leech-flow} satisfies
\begin{equation}
 \Lambda_s^*=\Lambda_{-s}\qquad(s\in\mathbb R).
 \label{eq:leech-dual}
\end{equation}
Its self-duality at a parameter \(s\) is equivalent to
\(D_{2s}\Lambda=\Lambda\).
\end{theorem}
\begin{proof}
For an invertible operator \(A\),
\[
 (AL)^*=A^{-\dagger}L^*.
\]
Indeed, \(\langle y,Ax\rangle_B=\langle A^\dagger y,x\rangle_B\)
is an integer for every \(x\in L\) precisely when
\(A^\dagger y\in L^*\). Taking \(A=D_s\), self-adjointness,
the identity \(D_s^{-1}=D_{-s}\), and \(\Lambda^*=\Lambda\) prove
\eqref{eq:leech-dual}. The equality
\(D_s\Lambda=D_{-s}\Lambda\), upon multiplication by \(D_s\), is equivalent
to \(D_{2s}\Lambda=\Lambda\).
\end{proof}

Covolume controls the volume of a fundamental domain;
self-duality controls all integer pairings with the lattice.
The first property is continuous along this collection. The
second imposes the exact arithmetic condition of the preceding theorem.
We shall therefore use the expression \emph{covolume one} for the
Euclidean flow and specify integrality whenever arithmetic
unimodularity is involved.

\begin{proposition}[Witness to local loss of integrality]
\label{prop:leech-nonintegral}
There exists \(\delta>0\) such that, if \(0<|s|<\delta\), the lattice
\(\Lambda_s\) is nonintegral and is not isometric to the Leech lattice.
\end{proposition}
\begin{proof}
The vector \(v_\alpha/3\) belongs to \(\Lambda\). Its transported
norm is
\begin{equation}
 f(s)=\left\|D_s\frac{v_\alpha}{3}\right\|_B^2
 =\frac29\left(16e^{2su_1}+
                       \sum_{j=2}^{12}e^{2su_j}\right).
 \label{eq:leech-witness}
\end{equation}
By \(\sum_j u_j=0\) and the first coordinate in
\eqref{eq:leech-u-coordinates},
\begin{equation}
\begin{split}
 f(0)&=6,\\
 f'(0)&=\frac49\left(16u_1+\sum_{j=2}^{12}u_j\right)
       =\frac{20}{3}u_1
       =-825-\frac{466}{3}\sqrt5\ne0.
\end{split}
\label{eq:leech-witness-derivative}
\end{equation}
Continuity of \(f'\) allows us to choose an interval
\((-\delta,\delta)\) on which \(f\) is strictly monotone.
Shrinking \(\delta\) also gives
\(11/2<f(s)<13/2\). For \(0<|s|<\delta\), we obtain \(f(s)\ne6\);
six is the only integer in that interval. Thus \(\Lambda_s\)
contains a vector of noninteger norm. An integral lattice has
integer norm on all its vectors. Since Leech is integral
and isometries preserve norms, the second statement is also
proved.
\end{proof}

The deformation preserves volume and ternary symmetry while
actually changing the arithmetic metric structure. Opposite
pairs of coordinates of \(u\) define a permutation of the ambient
space interchanging dilations and contractions. Promoting it
to an automorphism of the glued lattice would require preservation of the code
and the neighbour. The dual reciprocity already proved uses
only self-adjointness of the flow and initial self-duality.

\subsection{Convergence and reciprocity of theta functions}
\label{subsec:leech-theta}

For \(s\in\mathbb R\) and \(t>0\), define
\begin{equation}
 \Theta_s(t)=\sum_{\lambda\in\Lambda}
                \exp\bigl(-\pi t\|D_s\lambda\|_B^2\bigr).
 \label{eq:leech-theta-real}
\end{equation}
This is the Gaussian sum of all lengths of the deformed lattice,
with their multiplicities. If \(M=\max_j|u_j|\), we have
\begin{equation}
 e^{-2|s|M}\|x\|_B^2
 \leq\|D_sx\|_B^2
 \leq e^{2|s|M}\|x\|_B^2.
 \label{eq:leech-quadratic-bounds}
\end{equation}
On a compact subset of \(\mathbb R\times(0,\infty)\), the lower
bound provides a dominating Gaussian on the fixed lattice.
The series converges absolutely and uniformly. Derivatives of
finite order add polynomial factors in the coordinates of
\(\lambda\); those factors are dominated by another Gaussian.
Term-by-term differentiation is therefore
legitimate on those compact sets.

\begin{theorem}[Gaussian reciprocity]
\label{thm:leech-theta}
With the Euclidean volume associated with \(B\),
\begin{equation}
 \Theta_s(t)=t^{-12}\Theta_{-s}(t^{-1}),
 \qquad s\in\mathbb R,\quad t>0.
 \label{eq:leech-theta-reciprocity}
\end{equation}
In particular, \(\Theta_s(1)=\Theta_{-s}(1)\).
\end{theorem}
\begin{proof}
We fix the Fourier transform
\[
 \widehat f(y)=\int_E
      f(x)e^{-2\pi i\langle x,y\rangle_B}\,dx.
\]
For a lattice \(L\subset E\) and
\(f_t(x)=e^{-\pi t\|x\|_B^2}\), we periodize
\[
 F_t(x)=\sum_{\ell\in L}f_t(x+\ell).
\]
Normal convergence of the series and its derivatives gives a smooth
function on \(E/L\). If \(\xi\in L^*\), its Fourier coefficient is
\[
\begin{split}
 c_\xi
 &=\frac1{\operatorname{covol}(L)}
   \int_{E/L}F_t(x)e^{-2\pi i\langle x,\xi\rangle_B}\,dx\\
 &=\frac{\widehat f_t(\xi)}{\operatorname{covol}(L)}.
\end{split}
\]
The second equality decomposes \(E\) into translates of a fundamental
domain; the character has value one on \(L\). The product
of the one-dimensional Gaussian integrals, in a basis
orthonormal for \(B\), gives
\[
 \widehat f_t(\xi)=
 t^{-12}\exp\bigl(-\pi\|\xi\|_B^2/t\bigr).
\]
The Fourier series is absolutely convergent. Evaluating it at
\(x=0\), we obtain Poisson summation
\[
 \sum_{\ell\in L}e^{-\pi t\|\ell\|_B^2}
 =\frac{t^{-12}}{\operatorname{covol}(L)}
      \sum_{\xi\in L^*}e^{-\pi\|\xi\|_B^2/t}.
\]
Now \(L=\Lambda_s\) has covolume one and
\(L^*=\Lambda_{-s}\) by Theorem~\ref{thm:leech-dual}.
Substitution proves \eqref{eq:leech-theta-reciprocity}.
The case \(t=1\) is immediate.
\end{proof}

The preceding periodization is the Gaussian proof of Poisson summation;
it can be compared with \cite[Theorem 2, pp.~10--11]{Elkies}.
The positive-sign Fourier convention leads to the same
Gaussian transform, by its evenness. The theta equality
concerns the complete series and preserves the analytic domain of
convergence.

\begin{proposition}[Holomorphic extension]
\label{prop:leech-theta-complex}
For \(\tau\) in the upper half-plane
\(\mathbb H=\{\tau\in\mathbb C:\operatorname{Im}\tau>0\}\), the series
\begin{equation}
 \theta_s(\tau)=\sum_{\lambda\in\Lambda}
             e^{\pi i\tau\|D_s\lambda\|_B^2}
 \label{eq:leech-theta-complex}
\end{equation}
converges normally and satisfies
\begin{equation}
 \theta_s(-1/\tau)=(-i\tau)^{12}\theta_{-s}(\tau).
 \label{eq:leech-theta-modular}
\end{equation}
\end{proposition}
\begin{proof}
The absolute value of each summand is
\(e^{-\pi\operatorname{Im}\tau\|D_s\lambda\|_B^2}\).
On compact subsets of the upper half-plane, the imaginary part has
a positive lower bound, so
\eqref{eq:leech-quadratic-bounds} proves normal convergence and
holomorphy. For \(\tau=it\), with \(t>0\),
\eqref{eq:leech-theta-reciprocity} applied at \(t^{-1}\) gives
\(\theta_s(i/t)=t^{12}\theta_{-s}(it)\).
Both sides of \eqref{eq:leech-theta-modular} are holomorphic
on \(\mathbb H\); the identity theorem extends the equality
from the positive imaginary axis.
\end{proof}

Modular inversion relates the parameters \(s\) and \(-s\).
Additional periodicity under \(\tau\mapsto\tau+1\) uses
evenness of the norms. Proposition~\ref{prop:leech-nonintegral}
exhibits parameters arbitrarily close to zero for which
Euclidean integrality is lost. The collection therefore preserves
the proved analytic reciprocity, while a realization
through scalar modular forms or lattice vertex
algebras additionally incorporates the corresponding arithmetic
hypotheses. In the initial exceptional realization, those
hypotheses support continuation towards vertex algebras
and Monstrous Moonshine \cite{HMT-IV}.

\subsection{Split form and self-duality in dimension forty-eight}
\label{subsec:leech-split}

The loss of integrality of one Euclidean projection leads us to
consider both projections simultaneously. On the space
\(E\oplus E\), introduce the bilinear form
\begin{equation}
 \mathcal B((x,p),(x',p'))=
       \langle x,p'\rangle_B+\langle p,x'\rangle_B.
 \label{eq:leech-split-form}
\end{equation}
Its products pair the two coordinates. Define
\begin{equation}
\begin{split}
 \mathcal D_s&=D_s\oplus D_{-s},\\
 \Gamma_s&=\mathcal D_s(\Lambda\oplus\Lambda)\\
 &=\{(D_s\lambda,D_{-s}\mu):\lambda,\mu\in\Lambda\}.
\end{split}
\label{eq:leech-split-lattice}
\end{equation}

\begin{theorem}[Integral self-duality for the split form]
\label{thm:leech-split}
The form \(\mathcal B\) has signature \((24,24)\).
For every \(s\in\mathbb R\), \(\mathcal D_s\) is an isometry of
\(\mathcal B\), and \(\Gamma_s\) is an integral, even, self-dual lattice
with respect to this form.
\end{theorem}
\begin{proof}
Introducing
\[
 p_L=\frac{x+p}{\sqrt2},\qquad
 p_R=\frac{x-p}{\sqrt2},
\]
yields
\[
 \mathcal B((x,p),(x,p))=\|p_L\|_B^2-\|p_R\|_B^2.
\]
Since \(B\) is positive in dimension \(24\), the signature is
\((24,24)\). Self-adjointness and reciprocity of the blocks give
\[
 \langle D_sx,D_{-s}p'\rangle_B=\langle x,p'\rangle_B,
\]
and similarly for the second cross product. Hence
\(\mathcal B(\mathcal D_s\xi,\mathcal D_s\eta)=\mathcal B(\xi,\eta)\).

On \(\Gamma_0=\Lambda\oplus\Lambda\), all cross products
are integers, and
\[
 \mathcal B((\lambda,\mu),(\lambda,\mu))
   =2\langle\lambda,\mu\rangle_B\in2\mathbb Z.
\]
This proves integrality and evenness. If \((x,p)\) belongs to the dual
of \(\Gamma_0\) for \(\mathcal B\), pairing with
\((\lambda,0)\) and \((0,\mu)\) requires respectively
\(p,x\in\Lambda^*=\Lambda\). The dual is therefore \(\Gamma_0\).
The isometry \(\mathcal D_s\) transports these three properties
to \(\Gamma_s\).
\end{proof}

The metric specifies the content of the conclusion. The collection
\(\Gamma_s\) preserves integrality for the cross products
of \(\mathcal B\), even when \(\Lambda_s\) loses integrality
for \(B\). As abstract lattices with split form, the
\(\Gamma_s\) are isometric. Their two Euclidean projections,
considered with their respective positive metrics, vary with
the parameter.

\begin{proposition}[Exchange of reciprocal projections]
\label{prop:leech-split-involution}
The involution \(\mathcal J(x,p)=(p,x)\) satisfies
\[
 \mathcal J\Gamma_s=\Gamma_{-s},\qquad
 \mathcal J\mathcal D_s\mathcal J=\mathcal D_{-s}.
\]
In the coordinates \((p_L,p_R)\), it fixes \(p_L\) and reverses the sign
of \(p_R\).
\end{proposition}
\begin{proof}
The exchange turns
\((D_s\lambda,D_{-s}\mu)\) into
\((D_{-s}\mu,D_s\lambda)\), which ranges over \(\Gamma_{-s}\).
Conjugation of the two blocks gives the second identity.
Finally, \(x+p\) is preserved and \(x-p\) reverses sign.
\end{proof}

The involution expresses a duality of the complete lattice
construction. Any eventual physical representation preserves, in addition to
the split form, the operators and identification data
proper to that realization.

\subsection{Lattice energy and unitary equivalence}
\label{subsec:leech-energy}

The reciprocal-radius energy
\(m^2/R^2+w^2R^2\) interchanges its two terms under
\((m,w,R)\mapsto(w,m,R^{-1})\), with a unitary conjugation
preserving the operator domains \cite{HMT-IV}.
The realization on the twelve marked planes uses
\begin{equation}
 \mathcal E_s(\lambda,\mu)=
       \|D_s\lambda\|_B^2+\|D_{-s}\mu\|_B^2,
 \qquad \lambda,\mu\in\Lambda.
 \label{eq:leech-energy}
\end{equation}
The energy is positive and satisfies
\(\mathcal E_s(\lambda,\mu)=\mathcal E_{-s}(\mu,\lambda)\).

\begin{theorem}[Positive operator and conjugation of domains]
\label{thm:leech-energy}
On the Hilbert space \(\ell^2(\Lambda\oplus\Lambda)\), let
\[
 (H_s\psi)(\lambda,\mu)=
       \mathcal E_s(\lambda,\mu)\psi(\lambda,\mu)
\]
with domain
\begin{equation}
 \mathcal D(H_s)=
 \left\{\psi\in\ell^2(\Lambda\oplus\Lambda):
 \sum_{\lambda,\mu}\mathcal E_s(\lambda,\mu)^2
             |\psi(\lambda,\mu)|^2<\infty\right\}.
 \label{eq:leech-energy-domain}
\end{equation}
Then \(H_s\) is positive and self-adjoint with compact resolvent.
The map \(U\psi(\lambda,\mu)=\psi(\mu,\lambda)\) is unitary, and
\begin{equation}
 U\mathcal D(H_s)=\mathcal D(H_{-s}),\qquad
 UH_sU^{-1}=H_{-s}.
 \label{eq:leech-unitary}
\end{equation}
\end{theorem}
\begin{proof}
Multiplication by a real function on a discrete set,
with the maximal domain \eqref{eq:leech-energy-domain}, is
self-adjoint. This can be checked directly using the orthonormal
basis of singleton-supported vectors: membership
in the adjoint domain requires the sequence obtained by
multiplying by \(\mathcal E_s\) to be square summable, exactly
the same condition. Moreover,
\[
 \langle\psi,H_s\psi\rangle
   =\sum_{\lambda,\mu}\mathcal E_s(\lambda,\mu)
                         |\psi(\lambda,\mu)|^2\geq0.
\]
By \eqref{eq:leech-quadratic-bounds},
\[
 \mathcal E_s(\lambda,\mu)\geq
 e^{-2|s|M}\bigl(\|\lambda\|_B^2+\|\mu\|_B^2\bigr).
\]
A lattice has finitely many points in any bounded set;
hence every sublevel set of \(\mathcal E_s\) is finite. The diagonal
coefficients \((1+\mathcal E_s)^{-1}\) tend to zero outside
finite sets, so \((H_s+I)^{-1}\) is a norm limit
of finite-rank operators and is compact.

Exchanging the two indices preserves the norm of
\(\ell^2(\Lambda\oplus\Lambda)\), so \(U\) is unitary
and \(U^{-1}=U\). Equality of energies upon reversing \(s\) transforms
the domain condition of \(H_s\) into that of \(H_{-s}\), and
evaluation on each pair \((\lambda,\mu)\) proves
\eqref{eq:leech-unitary}.
\end{proof}

Three preservation laws with distinct domains are thus connected.
On the Euclidean planes, the flow maintains covolume and ternary
symmetry and relates each lattice to the dual at the opposite parameter.
In the doubled space, the cross products preserve an
integral, even, self-dual lattice. In the operator realization, the
reciprocal exchange preserves energy through a unitary
equivalence of operators together with their domains. The theta formulas
express the same reciprocity in the convergent sum of
lengths of the entire lattice.
